% =============================================================================
% Brief progress summary for supervisor review.
% Self-contained: compiles standalone (pdflatex supervisor_summary.tex), no
% external figure/bib dependencies. Numbers are taken verbatim from
% results/metrics/*.json and cross-checked by scripts/verify_thesis_numbers.py
% (31/31 pass at time of writing). Full write-up: reports/thesis_final.md.
% =============================================================================
\documentclass[11pt]{article}
\usepackage[margin=2.5cm]{geometry}
\usepackage{booktabs}
\usepackage{amsmath}
\usepackage{hyperref}
\usepackage{enumitem}
\usepackage{xcolor}
\usepackage{titlesec}

\titleformat{\section}{\large\bfseries}{\thesection}{0.6em}{}
\titleformat{\subsection}{\bfseries}{\thesubsection}{0.6em}{}

\title{\vspace{-1em}Memory-Efficient, Explainable Pneumonia Detection\\
\large Progress Summary for Supervisor Review}
\author{Sajan Khadka}
\date{\today}

\begin{document}
\maketitle
\vspace{-2em}

\section*{Research Question}
Can a lightweight convolutional network for pneumonia screening from chest
X-rays be trained, compressed, and evaluated in a way that is
\emph{deployment-honest} --- i.e.\ with leakage-checked data, correctly
measured memory footprint, statistically justified model selection, and
every reported number traceable to a reproducible artifact --- rather than
optimizing a single headline accuracy figure? The architectural contribution
is intentionally modest; the methodological contribution (evaluation
rigor, quantization behavior, and reproducibility tooling) is the core of
the thesis.

\section*{Method}
\begin{enumerate}[leftmargin=1.4em]
  \item \textbf{Data integrity.} Automated duplicate and train/validation
        leakage detection on the Kermany pediatric chest X-ray dataset
        (5{,}216 train / 624 held-out test images) removed 26 duplicate
        images and fixed 9 leaked train--validation pairs prior to training.
  \item \textbf{Model comparison.} Three backbones (EfficientNet-B0,
        ResNet-18, MobileNetV3-Small) trained under an identical protocol
        for a fair head-to-head comparison.
  \item \textbf{Pre-registered model selection.} EfficientNet-B0 was
        selected using a multi-objective criterion fixed in advance
        (\S3.9 of the full thesis) --- calibration error and false-positive
        rate at a fixed operating point --- rather than by post-hoc AUC
        ranking, since MobileNetV3-Small has marginally higher AUC but
        substantially worse calibration and false-positive count.
  \item \textbf{Post-training INT8 quantization.} Dynamic and static
        quantization compared; per-tensor weight quantization was found to
        collapse the model (AUC $\approx 0.40$), while per-channel
        quantization preserves performance.
  \item \textbf{Explainability.} Grad-CAM++ heatmaps generated for correct,
        false-positive, and false-negative predictions (qualitative only;
        not clinically validated).
  \item \textbf{External validation probe.} Zero-shot evaluation on the
        RSNA Pneumonia Detection Challenge dataset (adult patients, $n=12{,}024$)
        as an exploratory domain-shift check.
  \item \textbf{Statistical apparatus.} Bootstrap 95\% confidence intervals
        (1{,}000 resamples), DeLong test for AUC comparison, McNemar test
        for paired error comparison, and calibration metrics (ECE, Brier
        score, calibration slope/intercept).
  \item \textbf{Reproducibility.} \texttt{make reproduce} runs the full
        pipeline end-to-end; \texttt{scripts/verify\_thesis\_numbers.py}
        asserts every number in the final report matches
        \texttt{results/metrics/*.json} verbatim (currently 31/31 pass).
\end{enumerate}

\section*{Key Results}

\begin{table}[h]
\centering
\caption{Model comparison on the held-out Kermany test set ($n=624$)}
\begin{tabular}{lrrrrr}
\toprule
Model & Params (M) & AUC & Sensitivity & Specificity & Size (MB) \\
\midrule
\textbf{EfficientNet-B0 (selected)} & 4.34 & 0.9678 & 0.9667 & 0.8376 & 17.66 \\
ResNet-18 & 11.31 & 0.9594 & 0.9923 & 0.6923 & 45.32 \\
MobileNetV3-Small & 1.08 & 0.9743 & 0.9974 & 0.6709 & 4.44 \\
\bottomrule
\end{tabular}
\end{table}

\begin{table}[h]
\centering
\caption{EfficientNet-B0: effect of static INT8 quantization}
\begin{tabular}{lrrr}
\toprule
Metric & FP32 & Static INT8 & Change \\
\midrule
Size (MB) & 17.67 & 5.22 & $-70.5\%$ \\
Mean latency (ms) & 133.4 & 16.3 & $\sim 8\times$ faster \\
Inference-time RAM increase (MB, median of 5) & 1156.1 & 39.2 & $\sim 30\times$ lower \\
Test AUC & 0.9678 & 0.9444 & $-2.34$ pts \\
\bottomrule
\end{tabular}
\end{table}

\noindent Additional findings: EfficientNet-B0 vs.\ MobileNetV3-Small AUC
difference is not significant by DeLong's test ($p \approx 0.21$--$0.23$),
but the two models differ significantly on paired errors (McNemar,
$p = 0.0015$ / $0.0003$) and calibration (ECE 0.035 vs.\ 0.053--0.056),
supporting the selection criterion in point 3 above. A lazy streaming data
pipeline uses $\sim$10.5$\times$ less RAM than naive full-dataset loading
(322.4~MB vs.\ 3{,}378.1~MB, medians of 5 isolated runs each). Zero-shot
external validation on RSNA gives
AUC $= 0.8892$ (95\% CI [0.8825, 0.8954]), interpreted as expected domain
shift rather than a failure, and driven substantially by a calibration
shift (ECE 0.035 $\rightarrow$ 0.119).

\section*{Independent Verification}
Prior to this summary, all headline numbers were re-derived directly from
the raw JSON result files (not just taken from the written report), and the
project's own verification script was re-run and confirmed to pass 31/31.
The automated test suite (9 files, $\sim$40 tests) executes without
failures. The git history additionally documents three self-corrections: an
earlier, inflated 5-fold cross-validation figure; an earlier, incorrectly
measured (\texttt{tracemalloc}-based) memory metric; and, most recently, a
quantized-model-size and a memory-measurement fix described below ---
evidence of methodological self-checking rather than post-hoc narrative
construction.

\section*{Memory-Measurement Reliability (most recent fix)}
Two issues were found and corrected this cycle. \textbf{(1)} Quantized model
size was computed by re-serializing the model's state dict in memory rather
than measuring the file actually saved to disk; for the static-PTQ variant
(saved as a TorchScript archive) this understated size by $\sim$1.7\%
(5.13~MB $\to$ 5.22~MB, $-71\% \to -70.5\%$ compression). \textbf{(2)}
Peak-RSS memory measurement used a background thread polling every 10~ms,
which can miss a short-lived memory spike between polls --- repeating the
identical measurement showed swings from 0.2~MB to over 400~MB for the same
configuration. Replacing this with the OS kernel's own high-water-mark
counter (\texttt{getrusage().ru\_maxrss}), measuring each variant in a
fresh isolated process, and reporting the median of 5 runs (full samples
kept in \texttt{results/metrics/memory\_profile.json}) reduced that spread
to normal run-to-run noise ($\sim$2$\times$, not $\sim$2000$\times$). This
raised the reported inference-time memory benefit of static INT8 from
$\sim6\times$ to $\sim30\times$ --- a genuine, more-reliable measurement of
a narrower quantity (memory added by inference alone, excluding
model-build/quantization overhead), not a discretionary re-framing.

\section*{Limitations (acknowledged in the thesis)}
\begin{itemize}[leftmargin=1.4em]
  \item Single primary training dataset (pediatric); external validation is
        a single exploratory zero-shot probe, not a multi-site clinical
        validation.
  \item No $k$-fold cross-validation --- single stratified split with
        bootstrap confidence intervals.
  \item Efficiency measured via CPU proxy metrics; no physical edge-device
        (e.g.\ Raspberry Pi) benchmark.
  \item Grad-CAM++ explanations are qualitative; no radiologist review of
        localization quality.
  \item Not evaluated or positioned as a clinical device.
\end{itemize}

\section*{Assessment}
The work is internally consistent, its quantitative claims are independently
traceable to reproducible artifacts, and its limitations are stated rather
than obscured. The main open question for supervisory discussion is scope:
whether the single-dataset/single-split design is sufficient for the final
submission, or whether $k$-fold cross-validation and/or a second full
external dataset should be added before defense.

\vspace{1em}
\noindent\textit{Full write-up:} \texttt{reports/thesis\_final.md} \quad
\textit{Defense reference:} \texttt{reports/THESIS\_HANDBOOK.md} \quad
\textit{Reproduce all results:} \texttt{make reproduce}

\end{document}
